\section{Formulario de la unidad 2}
\begingroup
\footnotesize
\setlength{\parindent}{0pt}
\setlength{\abovedisplayskip}{3pt}
\setlength{\belowdisplayskip}{3pt}
\setlength{\abovedisplayshortskip}{3pt}
\setlength{\belowdisplayshortskip}{3pt}

\subsection*{Funciones vectoriales y cinemática}
Para $\vec r(t)=(f_1(t),\ldots,f_n(t))$, las operaciones se realizan por componentes:
\[
\operatorname{Dom}(\vec r)=\bigcap_{i=1}^n\operatorname{Dom}(f_i),\qquad
\lim_{t\to a}\vec r(t)=\left(\lim_{t\to a}f_i(t)\right)_{i=1}^n.
\]
\[
\vec r\,'(t)=(f_1'(t),\ldots,f_n'(t)),\qquad
\int_a^b\vec r(t)\,dt=\left(\int_a^b f_i(t)\,dt\right)_{i=1}^n.
\]
\[
\vec v=\vec r\,',\qquad v=\|\vec r\,'\|,\qquad
\vec a=\vec r\,'',\qquad
\vec r(t)=\vec r(t_0)+\int_{t_0}^t\vec v(u)\,du.
\]
Longitud de arco y triedro de Frenet, con $\vec r\,'\ne\vec0$:
\[
L=\int_a^b\|\vec r\,'(t)\|\,dt,\qquad
s(t)=\int_a^t\|\vec r\,'(u)\|\,du,\qquad
\hat T=\frac{\vec r\,'}{\|\vec r\,'\|}.
\]
\[
\kappa=\left\|\frac{d\hat T}{ds}\right\|
=\frac{\|\hat T'\|}{\|\vec r\,'\|}
=\frac{\|\vec r\,'\times\vec r\,''\|}{\|\vec r\,'\|^3},\qquad
\hat N=\frac{\hat T'}{\|\hat T'\|},\qquad \hat B=\hat T\times\hat N.
\]
$\hat N$ y $\hat B$ requieren $\kappa>0$.
\[
\vec a=a_T\hat T+a_N\hat N,\qquad
a_T=\frac{\vec r\,'\cdot\vec r\,''}{\|\vec r\,'\|}=\frac{dv}{dt},\qquad
a_N=\frac{\|\vec r\,'\times\vec r\,''\|}{\|\vec r\,'\|}=\kappa v^2.
\]

\subsection*{Funciones escalares, niveles y límites}
Para $f:D\subseteq\mathbb R^n\to\mathbb R$:
\[
\operatorname{Im}(f)=\{f(\mathbf x):\mathbf x\in D\},\qquad
\text{nivel }c:\ f(\mathbf x)=c,\qquad
\operatorname{Graf}(f)=\{(x,y,f(x,y)):(x,y)\in D\}.
\]
Dominio: denominador no nulo; radicando de raíz par no negativo;
argumento de logaritmo positivo.

\[
f\text{ continua en }\mathbf a\ \Longleftrightarrow\
\lim_{\mathbf x\to\mathbf a}f(\mathbf x)=f(\mathbf a).
\]
Cota polar para un límite en $(a,b)$:
\[
x=a+\rho\cos\theta,\quad y=b+\rho\sin\theta,\qquad
|f(a+\rho\cos\theta,b+\rho\sin\theta)-L|\leq M(\rho)\longrightarrow0.
\]
La cota debe valer para todo $\theta$ admisible. Dos trayectorias con límites
 distintos prueban que el límite no existe.

\subsection*{Derivadas parciales y regla de la cadena}
En un punto interior del dominio, si los límites existen:
\[
f_x(a,b)=\lim_{h\to0}\frac{f(a+h,b)-f(a,b)}h,\qquad
f_y(a,b)=\lim_{h\to0}\frac{f(a,b+h)-f(a,b)}h.
\]
\[
f_{xy}=\frac{\partial}{\partial y}\left(\frac{\partial f}{\partial x}\right),\qquad
f_{yx}=\frac{\partial}{\partial x}\left(\frac{\partial f}{\partial y}\right),\qquad
f\in C^2\ \Longrightarrow\ f_{xy}=f_{yx}.
\]
Para $f$ diferenciable y funciones interiores de clase $C^1$:
\begin{align*}
z=f(x(t),y(t)):&\qquad \frac{dz}{dt}=f_x\frac{dx}{dt}+f_y\frac{dy}{dt},\\
z=f(x(s,t),y(s,t)):&\qquad z_s=f_xx_s+f_yy_s,\qquad z_t=f_xx_t+f_yy_t,\\
w=f(x_1(\mathbf t),\ldots,x_m(\mathbf t)):&\qquad
\frac{\partial w}{\partial t_k}=\sum_{i=1}^m\frac{\partial f}{\partial x_i}\frac{\partial x_i}{\partial t_k}.
\end{align*}
Derivación implícita, cuando la ecuación define la función diferenciable indicada:
\begin{align*}
F(x,y(x))=0:&\quad y'=-\frac{F_x}{F_y}\quad(F_y\ne0),\\
F(x,y,z(x,y))=0:&\quad z_x=-\frac{F_x}{F_z},\quad z_y=-\frac{F_y}{F_z}\quad(F_z\ne0).
\end{align*}

\subsection*{Plano tangente, linealización y diferencial}
Para $f$ diferenciable en $(a,b)$, con $z_0=f(a,b)$:
\[
z-z_0=f_x(a,b)(x-a)+f_y(a,b)(y-b).
\]
\[
L(x,y)=f(a,b)+f_x(a,b)(x-a)+f_y(a,b)(y-b),\qquad f(x,y)\approx L(x,y).
\]
\[
dz=f_x\,dx+f_y\,dy,\qquad \Delta z\approx dz,\qquad
df=\sum_{i=1}^n\frac{\partial f}{\partial x_i}\,dx_i.
\]

\subsection*{Gradiente y derivada direccional}
Para $f$ diferenciable y $\|\mathbf u\|=1$:
\[
\nabla f=\left\langle\frac{\partial f}{\partial x_1},\ldots,
\frac{\partial f}{\partial x_n}\right\rangle,\qquad
D_{\mathbf u}f(\mathbf a)=\lim_{h\to0}\frac{f(\mathbf a+h\mathbf u)-f(\mathbf a)}h
=\nabla f(\mathbf a)\cdot\mathbf u.
\]
\[
\mathbf u=\frac{\mathbf v}{\|\mathbf v\|}\ (\mathbf v\ne\mathbf0),\qquad
\max_{\|\mathbf u\|=1}D_{\mathbf u}f=\|\nabla f\|,\qquad
\min_{\|\mathbf u\|=1}D_{\mathbf u}f=-\|\nabla f\|.
\]
Si $\nabla f\ne\mathbf0$, las direcciones extremas son
$\mathbf u_{\max}=\nabla f/\|\nabla f\|$ y
$\mathbf u_{\min}=-\nabla f/\|\nabla f\|$.
El gradiente es normal al nivel regular $f=c$.

\subsection*{Máximos, mínimos y multiplicadores de Lagrange}
Candidatos interiores: $\nabla f=\mathbf0$ y puntos donde alguna parcial no existe.
Para $f\in C^2$ cerca del punto crítico $(a,b)$:
\[
D=f_{xx}(a,b)f_{yy}(a,b)-[f_{xy}(a,b)]^2.
\]
\[
\begin{array}{c|c}
\text{Condición}&\text{Clasificación}\\\hline
D>0,\ f_{xx}(a,b)>0&\text{Mínimo local}\\
D>0,\ f_{xx}(a,b)<0&\text{Máximo local}\\
D<0&\text{Punto de silla}\\
D=0&\text{Criterio no concluyente}
\end{array}
\]
Para $f$ continua en una región no vacía, cerrada y acotada, los extremos
absolutos se obtienen comparando candidatos interiores y valores extremos
sobre toda la frontera, incluidos vértices y extremos de los arcos.

Para $f,g,h\in C^1$, los candidatos regulares con restricciones satisfacen
\[
\text{Una restricción:}\quad
\begin{cases}\nabla f=\lambda\nabla g,\\g=k,\end{cases}
\qquad \nabla g\ne\mathbf0.
\]
\[
\text{Dos restricciones:}\quad
\begin{cases}\nabla f=\lambda\nabla g+\mu\nabla h,\\g=k,\\h=m,\end{cases}
\qquad \nabla g,\nabla h\text{ linealmente independientes}.
\]
Se comparan los valores de $f$ en los candidatos; los puntos donde falla la
regularidad de las restricciones se analizan por separado.
\endgroup
